% c07 原书 p.75–90（PDF p.082–097）
% 续 §1.4 凝聚层；本块不含 \subsection。

设 $X$ 是环层空间，$\cF$ 是 $\cO_X$-模。定义 $\cF$ 的\emph{支集}\idx{zhiji2}{支集}为集合
\[
\operatorname{supp}\cF=\{P\in X\mid\cF_P\ne0\}.
\]

\begin{proposition}\label{it:1.4.14}
设 $X$ 是概形，$\cF$ 是拟凝聚 $\cO_X$-模。设 $\cF$ 是有限型的。则 $\operatorname{supp}\cF$ 是 $X$ 的闭子集。再设 $X$ 是诺特概形。则对 $\cO_X$ 的每个凝聚理想层 $\cI$，只要 $\operatorname{supp}\cF\subset\operatorname{supp}(\cO_X/\cI)$，就存在自然数 $n$ 使 $\cI^n\cF=0$。
\end{proposition}

\begin{proof}
用仿射开子概形 $U_i=\Spec A_i$（$i\in I$）覆盖 $X$，使 $\cF|_{U_i}\cong M_i^\sim$，其中 $M_i$ 是有限生成 $A_i$-模。由 1.3.28，有
\[
\operatorname{supp}\cF\cap U_i=V\bigl(\operatorname{Ann}_{A_i}(M_i)\bigr).
\]
故 $\operatorname{supp}\cF$ 是闭的。对每个 $i$，设 $I_i$ 是 $A_i$ 的理想，使 $\cI|_{U_i}\cong I_i^\sim$。则
\[
\operatorname{supp}(\cO_X/\cI)\cap U_i=V(I_i).
\]
若 $\operatorname{supp}\cF\subset\operatorname{supp}(\cO_X/\cI)$，则 $V(\operatorname{Ann}_{A_i}(M_i))\subset V(I_i)$。故
\[
I_i\subset\sqrt{\operatorname{Ann}_{A_i}(M_i)}.
\]
若 $X$ 是诺特概形，则 $I_i$ 有限生成，从而 $I_i^{n_i}\subset\operatorname{Ann}_{A_i}(M_i)$ 对某个自然数 $n_i$ 成立，即 $I_i^{n_i}M_i=0$。可取上述开覆盖为有限覆盖。故存在自然数 $n$ 使 $\cI^n\cF=0$。
\end{proof}

设 $(f,f^\sharp)\colon(Z,\cO_Z)\to(X,\cO_X)$ 是概形的闭浸入。态射 $f^\sharp\colon\cO_X\to f_*\cO_Z$ 的核 $\cI_Z$ 称为该闭浸入的\emph{理想层}\idx{lixiangceng}{理想层}。

\begin{proposition}\label{it:1.4.15}
设 $X$ 是概形。每个闭浸入到 $X$ 的理想层都拟凝聚。反之，对 $\cO_X$ 的每个拟凝聚理想层 $\cI$，环层空间 $(\operatorname{supp}(\cO_X/\cI),(\cO_X/\cI)|_{\operatorname{supp}(\cO_X/\cI)})$ 是概形，且由包含 $i\colon\operatorname{supp}(\cO_X/\cI)\to X$ 与典范态射
\[
i^\sharp\colon\cO_X\to i_*i^{-1}(\cO_X/\cI)=i_*\bigl((\cO_X/\cI)|_{\operatorname{supp}(\cO_X/\cI)}\bigr)
\]
组成的对 $(i,i^\sharp)$ 是以 $\cI$ 为理想层的闭浸入。此外，任何两个以 $\cI$ 为理想层的闭浸入到 $X$ 都同构。
\end{proposition}

\begin{proof}
闭浸入 $f\colon Z\to X$ 拟紧且分离。由 1.4.9 (iii)，$f_*\cO_Z$ 拟凝聚。故 $f^\sharp\colon\cO_X\to f_*\cO_Z$ 的核拟凝聚，即 $f$ 的理想层拟凝聚。

反之，给定拟凝聚理想层 $\cI$，记 $(Z,\cO_Z)=(\operatorname{supp}(\cO_X/\cI),(\cO_X/\cI)|_{\operatorname{supp}(\cO_X/\cI)})$。为证 $(Z,\cO_Z)$ 是概形，且 $(i,i^\sharp)\colon(Z,\cO_Z)\to(X,\cO_X)$ 是以 $\cI$ 为理想层的闭浸入，可设 $X=\Spec A$ 仿射。则 $\cI=I^\sim$，其中 $I$ 是 $A$ 的理想。有 $\cO_X/\cI=(A/I)^\sim$ 且 $\operatorname{supp}(\cO_X/\cI)=V(I)$。考虑由典范同态 $A\to A/I$ 诱导的闭浸入 $\varphi\colon\Spec A/I\to\Spec A$。它诱导同胚 $\Spec A/I\to V(I)$，且由 1.4.3 (i)，$\varphi_*\cO_{\Spec A/I}=(A/I)^\sim$。态射 $\varphi^\sharp\colon\cO_{\Spec A}\to\varphi_*\cO_{\Spec A/I}$ 的核是 $I^\sim$。故 $(Z,\cO_Z)$ 作为环层空间同构于 $\Spec A/I$，态射 $(i,i^\sharp)\colon(Z,\cO_Z)\to(X,\cO_X)$ 可等同于闭浸入 $\varphi$，且闭浸入 $i$ 的理想层是 $\cI$。

下面证明以 $\cI$ 为理想层的闭子概形的唯一性。设 $f\colon Z\to X$ 是以 $\cI$ 为理想层的闭浸入。$f^\sharp$ 诱导同构
\[
\cO_X/\cI\cong f_*\cO_Z.
\]
因 $f$ 把 $Z$ 同胚到 $X$ 的闭子集上，有 $\operatorname{supp}(f_*\cO_Z)=f(Z)$。故 $\operatorname{supp}(\cO_X/\cI)=f(Z)$，且 $f^\sharp$ 诱导同构 $f^{-1}(\cO_X/\cI)\cong\cO_Z$。因此 $Z$ 同构于闭子概形 $(\operatorname{supp}(\cO_X/\cI),(\cO_X/\cI)|_{\operatorname{supp}(\cO_X/\cI)})$。
\end{proof}

注意 1.4.15 给出 1.3.16 (ii) 的另一个证明。

\begin{proposition}\label{it:1.4.16}
设 $f\colon X\to Y$ 是概形的拟紧且拟分离态射。则存在闭浸入 $i\colon Z\to Y$ 与态射 $g\colon X\to Z$，使 $f=ig$，且 $i$ 满足如下泛性质：对每个分解 $f=i'g'$——其中 $i'\colon Z'\to Y$ 是闭浸入、$g'\colon Z'\to Y$ 是态射——存在唯一态射 $j\colon Z\to Z'$，使 $i'j=i$ 且 $jg=g'$。有 $i(Z)=\overline{f(X)}$。
\end{proposition}

\begin{tnote}
上面的泛性质中，原书把第二个因子写成 $g'\colon Z'\to Y$，但 $f=i'g'$ 要求 $g'$ 的靶是 $Z'$，即应为 $g'\colon X\to Z'$。这是原书的笔误，译文照原书直译，读者按 $g'\colon X\to Z'$ 理解即可。又，本命题陈述末尾原书作「The closed subscheme $Z$ the is called…」，多出一个 the，亦为笔误。[依据：原书 p.77 命题 1.4.16 及其证明]
\end{tnote}

闭子概形 $Z$ 称为 $f$ 的\emph{概形论像}\idx{gaixinglunxiang}{概形论像}。

\begin{proof}
设 $\cI$ 是 $f^\sharp\colon\cO_Y\to f_*\cO_X$ 的核。因 $f$ 拟紧且拟分离，由 1.4.9 (iii)，$f_*\cO_X$ 拟凝聚。故 $\cI$ 是 $\cO_Y$ 的拟凝聚理想层。设
\[
(Z,\cO_Z)=\bigl(\operatorname{supp}(\cO_Y/\cI),(\cO_Y/\cI)|_{\operatorname{supp}(\cO_Y/\cI)}\bigr)
\]
是由理想层 $\cI$ 定义的 $Y$ 的闭子概形，并设 $i\colon Z\to Y$ 是闭浸入。因 $f^\sharp$ 诱导单同态 $\cO_Y/\cI\rightarrowtail f_*\cO_X$，有
\[
Z=\operatorname{supp}(\cO_Y/\cI)=\operatorname{supp}(f_*\cO_X).
\]
这里用到了如下事实：若 $\cF$ 是 $Y$ 上的环层，则点 $P$ 属于 $\operatorname{supp}\cF$ 当且仅当 $0\ne1\in\cF_P$。有
\[
f(X)\subset\operatorname{supp}(f_*\cO_X)\subset\overline{f(X)}.
\]
因 $Z$ 是闭的，必须有 $Z=\overline{f(X)}$。

如下定义态射 $(g,g^\sharp)\colon(X,\cO_X)\to(Z,\cO_Z)$：底空间上的映射 $g\colon X\to Z$ 由 $f$ 诱导。态射 $f^\sharp\colon\cO_Y\to f_*\cO_X$ 诱导态射
\[
\cO_Y/\cI\to f_*\cO_X=i_*g_*\cO_X,
\]
从而诱导态射
\[
g^\sharp\colon\cO_Z=i^{-1}(\cO_Y/\cI)\to g_*\cO_X.
\]
可以验证 $f=ig$ 作为概形态射成立。

设 $i'\colon Z'\to Y$ 是闭浸入，且存在态射 $g'\colon X\to Z'$ 满足 $f=i'g'$。则 $f(X)\subset i'(Z')$，从而 $\overline{f(X)}\subset i'(Z')$。故存在连续映射 $j\colon Z\to Z'$，使 $i=i'j$ 作为拓扑空间之间的映射。设 $\cI'$ 是 $Z'$ 的理想层，即 $i'^\sharp\colon\cO_Y\to i'_*\cO_{Z'}$ 的核。因 $f=i'g'$，$\cI'$ 含于 $f^\sharp\colon\cO_Y\to f_*\cO_X$ 的核中，即 $\cI'$ 是 $\cI$ 的子层。故有典范态射 $\cO_Y/\cI'\to\cO_Y/\cI$。有 $\cO_Y/\cI'\cong i'_*\cO_{Z'}$ 与 $\cO_Y/\cI\cong i_*\cO_Z$。故有态射 $i'_*\cO_{Z'}\to i_*\cO_Z=i'_*j_*\cO_Z$，它诱导态射 $\cO_{Z'}\to j_*\cO_Z$。定义 $j^\sharp$ 为这最后一个态射。可以验证 $i'j=i$ 作为概形态射成立。于是有
\[
i'jg=ig=f=i'g'.
\]
因 $i'$ 是闭浸入，它是概形范畴中的单态射。故 $jg=g'$ 作为概形态射成立。使 $i'j=i$ 的 $j$ 的唯一性来自 $i'$ 是单态射这一事实。
\end{proof}

\begin{corollary}\label{it:1.4.17}
设 $f\colon Z\to X$ 是浸入。若 $f$ 拟紧，则 $f$ 可分解为 $Z\to F\to X$，使第一个箭头是开浸入，第二个箭头是闭浸入。
\end{corollary}

\begin{proof}
因 $f$ 是浸入，它分离。若 $f$ 拟紧，则由 1.4.16 可把 $f$ 分解为复合
\[
Z\xrightarrow{g}F\xrightarrow{i}X,
\]
使 $F$ 是 $f$ 的概形论像。则 $i\colon F\to X$ 是闭浸入。因 $i(F)=\overline{f(Z)}$ 且 $f$ 是浸入，$g\colon Z\to F$ 把 $Z$ 同胚到 $F$ 的开子集上。由概形论像的定义，对每个 $P\in Z$，同态
\[
g^\sharp_P\colon\cO_{F,g(P)}\to\cO_{Z,P}
\]
是单的。因 $f$ 是浸入，同态 $f^\sharp_P\colon\cO_{X,f(P)}\to\cO_{Z,P}$ 是满的。故 $g^\sharp_P$ 是满的。于是 $g^\sharp_P$ 是同构。由 1.3.15 (i)，$g$ 必是开浸入。
\end{proof}

\begin{proposition}\label{it:1.4.18}[Chow 引理]
设 $S$ 是诺特概形，$f\colon X\to S$ 是固有态射。则存在射影态射 $g\colon X'\to X$，使 $fg$ 射影，且 $g$ 对 $X$ 的某个开稠子集 $U$ 诱导同构 $g^{-1}(U)\xrightarrow{\sim}U$。
\end{proposition}

\begin{proof}
注意 $X$ 是诺特的。设 $X_i$（$i=1,\dots,n$）是 $X$ 的不可约分支。对每个 $i$，设 $V_i=X_i-\bigcup_{j\ne i}X_j$。则 $V_i$ 在 $X$ 中开、在 $X_i$ 中稠，且开浸入 $V_i\hookrightarrow X$ 拟紧。在 $X_i$ 上给一个概形结构，使它成为 $V_i\hookrightarrow X$ 的概形论像。则 $V_i$ 也是 $X_i$ 的开子概形。设 $f_i\colon X_i\to S$ 是 $f\colon X\to S$ 与闭浸入 $X_i\to X$ 的复合。则 $f_i$ 固有。设我们能对每个 $f_i$ 证明命题。设 $g_i\colon X_i'\to X_i$ 是射影态射，使 $f_ig_i$ 射影，并设 $U_i$ 是 $X_i$ 的开稠子集，使 $g_i$ 诱导同构 $g_i^{-1}(U_i)\xrightarrow{\sim}U_i$。令 $X'=\coprod_{i=1}^n X_i'$，$g\colon X'\to X$ 为 $g_i$ 诱导的态射。则 $g$ 与 $fg$ 都射影。取 $U=\bigcup_{i=1}^n(U_i\cap V_i)$。则 $U$ 在 $X$ 中稠，且 $g$ 诱导从 $g^{-1}(U)$ 到 $U$ 的同构。故为证命题，可设 $X$ 不可约。

用有限多个仿射开子概形 $V_i=\Spec A_i$（$i\in I$）覆盖 $S$，并用有限多个仿射开子概形 $U_{ij}=\Spec B_{ij}$（$j\in J_i$）覆盖每个 $f^{-1}(V_i)$，使每个 $B_{ij}$ 是有限生成 $A_i$-代数。设 $b_1,\dots,b_m$ 是 $B_{ij}$ 的一组有限生成元。则 $A_i$-代数同态
\[
A_i[x_1,\dots,x_m]\to B_{ij},\quad x_k\mapsto b_k
\]
是满的。它诱导闭浸入 $U_{ij}\to\Spec A_i[x_1,\dots,x_m]$。另一方面，有同构
\[
A_i[x_1,\dots,x_m]\to(A_i[x_0,\dots,x_m])_{(x_0)},\quad x_k\mapsto\frac{x_k}{x_0},
\]
以及开浸入
\[
\Spec (A_i[x_0,\dots,x_m])_{(x_0)}\to\Proj A_i[x_0,\dots,x_m].
\]
故有浸入
\[
U_{ij}\to\Proj A_i[x_0,\dots,x_m]=\mathbb{P}_{A_i}^m.
\]
开浸入 $\Spec A_i=V_i\to S$ 诱导开浸入 $\mathbb{P}_{A_i}^m\to\mathbb{P}_S^m$。故最终得到浸入 $U_{ij}\to\mathbb{P}_S^m$。因 $U_{ij}$ 是诺特的，该浸入拟紧。由 1.4.17，可把此浸入分解为复合 $U_{ij}\to P_{ij}\to\mathbb{P}_S^m$，使第一个箭头是开浸入、第二个箭头是闭浸入。改变记号，可设 $X$ 可被非空开子集 $U_i$（$i=1,\dots,n$）覆盖，使每个态射 $U_i\to S$ 可分解为
\[
U_i\to P_i\to S,
\]
其中 $U_i\to P_i$ 是开浸入，$P_i\to S$ 射影。

设 $U=\bigcap_{i=1}^n U_i$。因 $X$ 不可约，$U$ 在 $X$ 中稠。考虑由开浸入 $U\to X$ 与态射 $U\to P_i$ 定义的态射
\[
h\colon U\to X\times_S P_1\times_S\cdots\times_S P_n.
\]
设 $P=P_1\times_S\cdots\times_S P_n$，设 $X'$ 是 $h\colon U\to X\times_S P$ 的概形论像，并设 $g\colon X'\to X$ 与 $g'\colon X'\to P$ 分别是闭浸入 $X'\to X\times_S P$ 与 $X\times_S P$ 到各因子投影的复合。

注意 $g\colon X'\to X$ 射影，因为它是闭浸入 $X'\to X\times_S P$ 与射影态射 $X\times_S P\to X$ 的复合。有 $g^{-1}(U)=X'\cap(U\times_S P)$。注意 $X'\cap(U\times_S P)\to U\times_S P$ 是 $h\colon U\to U\times_S P$ 的概形论像。但 $U\to U\times_S P$ 是闭浸入，因为它是 $U\to P$ 的图。故 $U\to U\times_S P$ 是它自身的概形论像。于是态射 $U\to X'$ 诱导同构 $U\cong X'\cap(U\times_S P)$。把它与 $X'\cap(U\times_S P)=g^{-1}(U)\xrightarrow{g}U$ 复合，得到 $\id_U$。故 $g$ 诱导从 $g^{-1}(U)$ 到 $U$ 的同构。

显然 $X'\cap(U_i\times_S P)$（$i=1,\dots,n$）构成 $X'$ 的开覆盖。设 $W_i$ 是 $U_i$ 在投影 $P\to P_i$ 下的原像。我们断言
\[
X'\cap(U_i\times_S P)\subset X'\cap(X\times_S W_i).
\]
这将推出 $X'\cap(X\times_S W_i)$（$i=1,\dots,n$）构成 $X'$ 的开覆盖。对每个 $i$，有交换图
\begin{center}
\begin{tikzcd}[column sep=3em,row sep=2.4em]
U \arrow[r] \arrow[d] & P \arrow[d] \\
U_i \arrow[r] & P_i,
\end{tikzcd}
\end{center}
其中上方水平箭头是态射 $U\to P$，它与每个投影 $P\to P_i$ 的复合等于复合 $U\hookrightarrow U_i\to P_i$。因 $U$ 在 $X'$ 中的像是稠的，下图作为底空间映射的图交换：
\begin{center}
\begin{tikzcd}[column sep=3em,row sep=2.4em]
X'\cap(U_i\times_S P) \arrow[r] \arrow[d] & P \arrow[d] \\
U_i \arrow[r] & P_i.
\end{tikzcd}
\end{center}
（对照习题 2。）这说明 $X'\cap(U_i\times_S P)\to P$ 的像含于 $W_i$。故 $X'\cap(U_i\times_S P)\subset X'\cap(X\times_S W_i)$。

最后证明 $g'\colon X'\to P$ 是闭浸入。这将推出 $X'$ 在 $S$ 上射影，从而 $fg\colon X'\to S$ 射影。因 $X$ 与 $P$ 在 $S$ 上固有，$X\times_S P$ 在 $S$ 上固有。因 $X'$ 是 $X\times_S P$ 的闭子概形，$X'$ 也在 $S$ 上固有。由 1.3.26 (vi)，$g'\colon X'\to P$ 固有。为证它是闭浸入，只需证它是浸入。有 $g'^{-1}(W_i)=X'\cap(X\times_S W_i)$。因 $X'\cap(X\times_S W_i)$（$i=1,\dots,n$）构成 $X'$ 的开覆盖，只需证每个 $g'\colon X'\cap(X\times_S W_i)\to W_i$ 是闭浸入。

考虑由复合 $W_i\to U_i\to X$ 定义的态射，其中第一个箭头由投影 $P\to P_i$ 诱导。设 $\Gamma\colon W_i\to X\times_S W_i$ 是此态射的图。它是闭浸入。设 $U\to P$ 是使与每个投影 $P\to P_i$ 的复合等于 $U\hookrightarrow U_i\to P_i$ 的态射。$U\to P$ 的像含于 $W_i$，故它诱导态射 $U\to W_i$，且态射 $U\to X\times_S P$ 诱导态射 $U\to X\times_S W_i$。此外有交换图
\begin{center}
\begin{tikzcd}[column sep=3.2em,row sep=2.6em]
& W_i \arrow[d,"\Gamma"] \\
U \arrow[ur] \arrow[r] & X\times_S W_i.
\end{tikzcd}
\end{center}
因 $X'\cap(X\times_S W_i)\to X\times_S W_i$ 是 $U\to X\times_S W_i$ 的概形论像，态射 $X'\cap(X\times_S W_i)\to X\times_S W_i$ 通过闭浸入 $W_i\to X\times_S W_i$ 分解：
\begin{center}
\begin{tikzcd}[column sep=3.2em,row sep=2.6em]
& W_i \arrow[d,"\Gamma"] \\
X'\cap(X\times_S W_i) \arrow[ur,"\theta"] \arrow[r] & X\times_S W_i.
\end{tikzcd}
\end{center}
注意 $\theta\colon X'\cap(X\times_S W_i)\to W_i$ 必是闭浸入。考虑交换图
\begin{center}
\begin{tikzcd}[column sep=4em,row sep=2.6em]
& W_i \arrow[d,"\Gamma"] \\
X'\cap(X\times_S W_i) \arrow[ur,"\theta"] \arrow[r] \arrow[dr,"g'"'] & X\times_S W_i \arrow[d] \\
& W_i.
\end{tikzcd}
\end{center}
因图中竖直箭头的复合是恒等，态射 $g'\colon X'\cap(X\times_S W_i)\to W_i$ 与 $\theta$ 相同，故 $g'$ 是闭浸入。这证明了断言。
\end{proof}

\begin{tnote}
Chow 引理把固有态射「化归为射影」：存在射影态射 $g\colon X'\to X$，使 $fg$ 射影，且 $g$ 在 $X$ 的某个稠密开子集上是同构。它在证明中把固有情形逐步约化到射影情形，是本节后续论证（以及层上同调）的常用工具。[依据：原书 p.78 命题 1.4.18 及 p.78–81 的证明]
\end{tnote}

设 $S=\bigoplus_{d=0}^\infty S_d$ 是分次环。\emph{分次 $S$-模}\idx{fencihuan}{分次环}是指 $S$-模 $M$，连同阿贝尔群的分解 $M=\bigoplus_{i=-\infty}^\infty M_i$，使 $S_dM_i\subset M_{d+i}$。$M_i$ 中的元素称为\emph{次数为 $i$ 的齐次元素}。对每个整数 $n$，定义 $M(n)$ 为如下分次 $S$-模：其下层的 $S$-模是 $M$，其分次由 $(M(n))_d=M_{d+n}$ 定义。对 $S$ 的每个齐次素理想 $\fp$，定义 $M_{(\fp)}$ 为 $S_{(\fp)}$-模
\[
M_{(\fp)}=\Bigl\{\frac{m}{t}\in M_{\fp}\Bigm| m\in M,\ t\in S-\fp \text{ 是次数相同的齐次元素}\Bigr\}.
\]
对 $S$ 的每个齐次元素 $f$，定义 $M_{(f)}$ 为 $S_{(f)}$-模
\[
M_{(f)}=\Bigl\{\frac{m}{f^k}\in M_f\Bigm| k\in\bN,\ m\in M \text{ 是次数为 } k\deg f \text{ 的齐次元素}\Bigr\}.
\]
如下定义 $\cO_{\Proj S}$-模层 $M^\sim$：对 $\Proj S$ 的每个开子集 $U$，$M^\sim(U)$ 由满足以下两条件的映射 $s\colon U\to\coprod_{\fp\in\Proj S}M_{(\fp)}$ 组成：
\begin{enumerate}[label=(\alph*),leftmargin=2.8em]
  \item 对每个 $\fp\in U$，有 $s(\fp)\in M_{(\fp)}$。
  \item 对每个 $\fp\in U$，存在 $\fp$ 的邻域 $U_{\fp}$，含于 $U$，以及次数相同的齐次元素 $m\in M$ 与 $f\in S$，使对每个 $\fq\in U_{\fp}$ 有 $f\notin\fq$ 且 $s(\fq)=\frac{m}{f}$ 在 $M_{(\fq)}$ 中。
\end{enumerate}
对开子集的每个包含 $V\subset U$，定义 $M^\sim(U)\to M^\sim(V)$ 为映射的限制。

\begin{proposition}\label{it:1.4.19}
设 $S$ 是分次环，$M$ 是分次 $S$-模。
\begin{enumerate}[label=(\roman*),leftmargin=2.8em]
  \item 对每个 $\fp\in\Proj S$，有 $(M^\sim)_{\fp}\cong M_{(\fp)}$。
  \item 对 $S$ 中每个正次数的齐次元素 $f$，$M^\sim|_{D_+(f)}$ 可经由同构 $D_+(f)\cong\Spec S_{(f)}$ 与 $(M_{(f)})^\sim$ 等同。
  \item $M^\sim$ 拟凝聚。若 $S$ 是诺特环且 $M$ 是有限生成 $S$-模，则 $\Proj S$ 是诺特概形，且 $M^\sim$ 凝聚。
\end{enumerate}
\end{proposition}

\noindent 留给读者证明此命题。我们只解释如下事实：若 $S$ 是诺特环且 $M$ 有限生成，则对 $S$ 的任意齐次元素 $f$，$M_{(f)}$ 是诺特的。设 $N_1\subset N_2\subset\cdots$ 是 $M_{(f)}$ 的子模的升链。对每个 $i$，设 $S_fN_i$ 是 $M_f$ 的由 $N_i$ 生成的子模。则有 $S_fN_1\subset S_fN_2\subset\cdots$。因 $M_f$ 是诺特的，当 $n$ 充分大时有 $S_fN_n=S_fN_{n+1}=\cdots$。又 $S_fN_i\cap M_{(f)}=N_i$。故 $N_n=N_{n+1}=\cdots$。于是 $M_{(f)}$ 是诺特的。

对任意整数 $n$，定义 $\cO_{\Proj S}(n)=S(n)^\sim$。称 $\cO_{\Proj S}(1)$ 为\emph{扭层}\idx{niuceng}{扭层}。对任意 $\cO_{\Proj S}$-模 $\cF$，定义
\[
\cF(n)=\cF\otimes_{\cO_{\Proj S}}\cO_{\Proj S}(n).
\]

设 $M$ 与 $N$ 是分次 $S$-模。如下定义分次 $S$-模 $M\otimes_SN$：有
\[
M\otimes_{\bZ}N=\Bigl(\bigoplus_dM_d\Bigr)\otimes_{\bZ}\Bigl(\bigoplus_eN_e\Bigr)=\bigoplus_n\Bigl(\bigoplus_{d+e=n}M_d\otimes_{\bZ}N_e\Bigr).
\]
它给出 $M\otimes_{\bZ}N$ 的分解。$M\otimes_SN$ 是 $M\otimes_{\bZ}N$ 对由形如 $sm\otimes n-m\otimes sn$ 的元素生成的子群的商。注意此子群是 $M\otimes_{\bZ}N$ 的齐次子群。故 $M\otimes_{\bZ}N$ 的分解诱导 $M\otimes_SN$ 的分次。

\begin{proposition}\label{it:1.4.20}
设 $S$ 是由 $S_1$ 作为 $S_0$-代数生成的分次环。
\begin{enumerate}[label=(\roman*),leftmargin=2.8em]
  \item 对每个 $n$，$\cO_{\Proj S}(n)$ 是可逆 $\cO_{\Proj S}$-模。
  \item 设 $M$ 与 $N$ 是分次 $S$-模。则
  \[
  (M\otimes_SN)^\sim\cong M^\sim\otimes_{\cO_{\Proj S}}N^\sim.
  \]
  特别地，有
  \[
  (M(n))^\sim=(M\otimes_SS(n))^\sim=M^\sim\otimes_{\cO_{\Proj S}}S(n)^\sim=M^\sim(n),
  \]
  \[
  (S(m+n))^\sim=(S(m)\otimes_SS(n))^\sim=S(m)^\sim\otimes_{\cO_{\Proj S}}S(n)^\sim.
  \]
  于是
  \[
  \cO_{\Proj S}(m+n)=\cO_{\Proj S}(m)\otimes_{\cO_{\Proj S}}\cO_{\Proj S}(n),
  \]
  \[
  \cO_{\Proj S}(n)=\cO_{\Proj S}(1)^{\otimes n},
  \]
  \[
  \cO_{\Proj S}(-1)=\mathcal{H}om_{\cO_{\Proj S}}\bigl(\cO_{\Proj S}(1),\cO_{\Proj S}\bigr).
  \]
  \item 设 $T$ 是另一个分次环，由 $T_1$ 作为 $T_0$-代数生成；设 $\phi\colon S\to T$ 是同态，且对所有 $d$ 满足 $\phi(S_d)\subset T_d$；并设
  \[
  U=\{\fq\in\Proj T\mid S_+\not\subset\phi^{-1}(\fq)\}.
  \]
  则 $U$ 是 $\Proj T$ 的开子集，且 $\phi$ 典范地诱导概形态射 $f\colon U\to\Proj S$。对每个分次 $S$-模 $M$，有
  \[
  f^*M^\sim=(T\otimes_SM)^\sim|_U,
  \]
  且对每个分次 $T$-模 $N$，有
  \[
  f_*(N^\sim|_U)=({}_SN)^\sim,
  \]
  其中右边把 $N$ 视为 $S$-模。特别地，有
  \[
  f^*\cO_{\Proj S}(n)\cong\cO_{\Proj T}(n)|_U.
  \]
\end{enumerate}
\end{proposition}

\begin{proof}
(i) 对每个齐次元素 $f\in S_+$，在 $D_+(f)\cong\Spec S_{(f)}$ 上有 $\cO_{\Proj S}(n)|_{D_+(f)}\cong(S(n)_{(f)})^\sim$。当 $f\in S_1$ 时，有 $S_{(f)}$-模同构
\[
S_{(f)}\to S(n)_{(f)},\quad x\mapsto f^nx.
\]
故 $\cO_{\Proj S}|_{D_+(f)}\cong\cO_{\Proj S}(n)|_{D_+(f)}$。因 $S$ 由 $S_1$ 作为 $S_0$-代数生成，$D_+(f)$（$f\in S_1$）构成 $\Proj S$ 的开覆盖。故 $\cO_{\Proj S}(n)$ 可逆。

(ii) 定义态射
\[
\phi\colon M^\sim\otimes_{\cO_{\Proj S}}N^\sim\to(M\otimes_SN)^\sim
\]
如下。因 $M^\sim\otimes_{\cO_{\Proj S}}N^\sim$ 是前层
\[
U\mapsto M^\sim(U)\otimes_{\cO_{\Proj S}(U)}N^\sim(U)
\]
相伴的层，只需定义从此前层到 $(M\otimes_SN)^\sim$ 的态射。对 $M^\sim(U)$ 的截面 $s\colon U\to\coprod M_{(\fp)}$ 与 $N^\sim(U)$ 的截面 $t\colon U\to\coprod N_{(\fp)}$，定义 $(M\otimes_SN)^\sim(U)$ 的截面
\[
\phi(s\otimes t)\colon U\to\coprod_{\fp\in\Proj S}(M\otimes_SN)_{(\fp)}
\]
使对每个 $\fp\in U$，$\phi(s\otimes t)(\fp)$ 是 $s(\fp)\otimes t(\fp)$ 在典范同态
\[
M_{(\fp)}\otimes_{S_{(\fp)}}N_{(\fp)}\to(M\otimes_SN)_{(\fp)}
\]
下的像。

对每个 $f\in S_1$，经由同构 $D_+(f)\cong\Spec S_{(f)}$，有
\[
(M^\sim\otimes_{\cO_{\Proj S}}N^\sim)|_{D_+(f)}\cong M_{(f)}^\sim\otimes_{\cO_{\Spec S_{(f)}}}N_{(f)}^\sim\cong\bigl(M_{(f)}\otimes_{S_{(f)}}N_{(f)}\bigr)^\sim,
\]
\[
(M\otimes_SN)^\sim|_{D_+(f)}\cong(M\otimes_SN)_{(f)}^\sim.
\]
因 $D_+(f)$（$f\in S_1$）构成 $\Proj S$ 的开覆盖，为证 $\phi$ 是同构，只需证同态
\[
\phi_{(f)}\colon M_{(f)}\otimes_{S_{(f)}}N_{(f)}\to(M\otimes_SN)_{(f)}
\]
对每个 $f\in S_1$ 是同构。显然它是满的。有交换图
\begin{center}
\begin{tikzcd}[column sep=3.2em,row sep=2.6em]
M_{(f)}\otimes_{S_{(f)}}N_{(f)} \arrow[r,"\phi_{(f)}"] \arrow[d] & (M\otimes_SN)_{(f)} \arrow[d] \\
M_f\otimes_{S_f}N_f \arrow[r,"\cong"'] & (M\otimes_SN)_f.
\end{tikzcd}
\end{center}
为证 $\phi_{(f)}$ 单，只需证第一个竖直箭头单。事实上它有左逆
\[
M_f\otimes_{S_f}N_f\to M_{(f)}\otimes_{S_{(f)}}N_{(f)},\qquad
\frac{\sum_im_i}{f^k}\otimes\frac{\sum_in_i}{f^l}\mapsto\Bigl(\sum_i\frac{m_i}{f^i}\Bigr)\otimes\Bigl(\sum_i\frac{n_i}{f^i}\Bigr),
\]
其中对每个 $i$，$m_i$ 与 $n_i$ 分别是 $M$ 与 $N$ 中次数为 $i$ 的齐次元素。留给读者验证此映射良定义。

(iii) 有 $U=\Proj T-V(\phi(S_+)T)$。故 $U$ 是开的。考虑映射
\[
f\colon U\to\Proj S,\quad\fq\mapsto\phi^{-1}(\fq).
\]
对 $S$ 的每个齐次理想 $\fa$，有 $f^{-1}(V(\fa))=V(\phi(\fa)T)\cap U$。故 $f$ 连续。如下定义 $f^\sharp\colon\cO_{\Proj S}\to f_*(\cO_{\Proj T}|_U)$：对 $\Proj S$ 的每个开子集 $V$ 与每个截面 $s\colon V\to\coprod_{\fp\in\Proj S}S_{(\fp)}$，定义截面
\[
f^\sharp(s)\colon f^{-1}(V)\to\coprod_{\fq\in\Proj T}T_{(\fq)}
\]
为
\[
(f^\sharp(s))(\fq)=\phi_\fq\bigl(s(f(\fq))\bigr)
\]
对每个 $\fq\in f^{-1}(V)$，其中 $\phi_\fq\colon S_{(\phi^{-1}(\fq))}\to T_{(\fq)}$ 是 $\phi$ 诱导的同态。可以验证 $(f,f^\sharp)$ 是概形态射。注意对每个齐次元素 $a\in S_+$，有 $f^{-1}(D_+(a))=D_+(\phi(a))$。

对每个 $T$-模 $N$，定义
\[
\alpha\colon({}_SN)^\sim\to f_*(N^\sim|_U)
\]
如下：对每个截面 $s\colon V\to\coprod_{\fp\in\Proj S}({}_SN)_{(\fp)}$，设
\[
\alpha(s)\colon f^{-1}(V)\to\coprod_{\fq\in\Proj T}N_{(\fq)}
\]
是 $f_*(N^\sim|_U)(V)=N^\sim(f^{-1}(V))$ 的截面，使对每个 $\fq\in f^{-1}(V)$，$\alpha(s)(\fq)$ 是 $s(\phi^{-1}(\fq))$ 在典范同态
\[
({}_SN)_{(\phi^{-1}(\fq))}\to N_{(\fq)},\quad\frac{n}{s}\mapsto\frac{n}{\phi(s)}
\]
下的像。

为证 $\alpha$ 是同构，只需证对每个齐次元素 $a\in S_+$，同态
\[
({}_SN)^\sim(D_+(a))\to f_*(N^\sim|_U)(D_+(a))
\]
是同构。有
\[
({}_SN)^\sim(D_+(a))\cong({}_SN)_{(a)},\qquad
f_*(N^\sim|_U)(D_+(a))\cong N^\sim(D_+(\phi(a)))\cong N_{(\phi(a))}.
\]
断言来自典范同态
\[
({}_SN)_{(a)}\to N_{(\phi(a))},\quad\frac{n}{a^k}\mapsto\frac{n}{\phi(a)^k}
\]
是同构。

对任意 $S$-模 $M$，$S$-模同态
\[
M\to{}_S(T\otimes_SM),\quad m\mapsto1\otimes m
\]
诱导态射
\[
M^\sim\to({}_S(T\otimes_SM))^\sim\cong f_*\bigl((T\otimes_SM)^\sim|_U\bigr).
\]
因 $f^*$ 左伴随于 $f_*$，得到态射
\[
\beta\colon f^*M^\sim\to(T\otimes_SM)^\sim|_U.
\]
对任意齐次元素 $a\in S_+$，有
\[
(f^*M^\sim)|_{D_+(\phi(a))}\cong(f|_{D_+(\phi(a))})^*\bigl(M^\sim|_{D_+(a)}\bigr)\cong(f|_{\Spec T_{(\phi(a))}})^*\bigl(M^\sim_{(a)}\bigr)\cong\bigl(T_{(\phi(a))}\otimes_{S_{(a)}}M_{(a)}\bigr)^\sim,
\]
\[
(T\otimes_SM)^\sim|_{D_+(\phi(a))}\cong(T\otimes_SM)_{(\phi(a))}.
\]
可以验证对任意 $a\in S_1$，$\beta$ 诱导同构
\[
T_{(\phi(a))}\otimes_{S_{(a)}}M_{(a)}\cong(T\otimes_SM)_{(\phi(a))}.
\]
对照 1.4.20 (ii) 的证明。因 $D_+(\phi(a))$（$a\in S_1$）构成 $U$ 的开覆盖，$\beta$ 是同构。
\end{proof}

留给读者证明如下引理：（对照 1.3.9 (iii) 与 1.4.4 的证明。）

\begin{lemma}\label{it:1.4.21}
设 $X$ 是概形，$\cL$ 是可逆 $\cO_X$-模，$f\in\cL(X)$。定义 $X_f$ 为 $X$ 的如下点 $x$ 组成的子集：$f$ 在 $x$ 处的芽不在 $\fm_x\cL_x$ 中，其中 $\fm_x$ 是局部环 $\cO_{X,x}$ 的极大理想。则 $X_f$ 是开的。设 $\cF$ 是拟凝聚 $\cO_X$-模。
\begin{enumerate}[label=(\roman*),leftmargin=2.8em]
  \item 设 $X$ 拟紧。若 $s\in\cF(X)$ 是限制到 $X_f$ 上为零的截面，则存在自然数 $n$，使截面 $s\otimes f^{\otimes n}\in(\cF\otimes_{\cO_X}\cL^{\otimes n})(X)$ 为零。
  \item 设 $X$ 拟紧且拟分离。给定截面 $t\in\cF(X_f)$，存在自然数 $n$，使截面 $t\otimes f^{\otimes n}\in(\cF\otimes_{\cO_X}\cL^{\otimes n})(X_f)$ 可延拓为 $(\cF\otimes_{\cO_X}\cL^{\otimes n})(X)$ 的截面。
\end{enumerate}
\end{lemma}

设 $S$ 是由 $S_1$ 作为 $S_0$-代数生成的分次环，$X=\Proj S$。记 $\cO_{\Proj S}(n)$ 为 $\cO(n)$。对任意 $\cO_X$-模 $\cF$，定义与 $\cF$ 相伴的分次 $S$-模为
\[
\Gamma_*(\cF)=\bigoplus_{n=-\infty}^\infty\cF(n)(X),
\]
其中 $S$ 在 $\Gamma_*(\cF)$ 上的纯量积由如下典范同态的复合定义：
\[
S_d\otimes_{\bZ}\cF(n)(X)\to\cO(d)(X)\otimes_{\bZ}\cF(n)(X)\to\cF(d+n)(X).
\]

\begin{proposition}\label{it:1.4.22}
设 $S$ 是由 $S_1$ 作为 $S_0$-代数有限生成的分次环，$\cF$ 是拟凝聚 $\cO_{\Proj S}$-模。则有典范同构 $\Gamma_*(\cF)^\sim\cong\cF$。
\end{proposition}

\begin{proof}
记 $\cO_{\Proj S}(n)$ 为 $\cO(n)$。对 $S$ 的每个次数为 $d$ 的齐次元素 $f$，定义
\[
\phi(D_+(f))\colon\Gamma_*(\cF)^\sim(D_+(f))\cong\Gamma_*(\cF)_{(f)}\to\cF(D_+(f))
\]
如下：对每个 $\frac{s}{f^k}\in\Gamma_*(\cF)_{(f)}$，其中 $s\in\cF(dk)(X)$，定义 $\phi(D_+(f))(\frac{s}{f^k})$ 为 $\frac{1}{f^k}\otimes s|_{D_+(f)}$ 在典范同态
\[
\cO(-dk)(D_+(f))\otimes_{\cO(D_+(f))}\cF(dk)(D_+(f))\to\cF(D_+(f))
\]
下的像。对 $\cL=\cO(d)$ 应用 1.4.21，可知 $\phi(D_+(f))$ 是同构。对 $X$ 的每个开子集 $U$，选 $U$ 的开覆盖 $\{D_+(f_i)\}_{i\in I}$，并定义
\[
\phi(U)\colon\Gamma_*(\cF)^\sim(U)\to\cF(U)
\]
为使下图交换的唯一同构：
\begin{center}
\begin{tikzcd}[column sep=1.2em,row sep=2.4em]
0 \arrow[r] & \Gamma_*(\cF)^\sim(U) \arrow[r] \arrow[d,"\phi(U)"] & \prod_{i\in I}\Gamma_*(\cF)^\sim(D_+(f_i)) \arrow[r] \arrow[d] & \prod_{i,j\in I}\Gamma_*(\cF)^\sim(D_+(f_i)\cap D_+(f_j)) \arrow[d] \\
0 \arrow[r] & \cF(U) \arrow[r] & \prod_{i\in I}\cF(D_+(f_i)) \arrow[r] & \prod_{i,j\in I}\cF(D_+(f_i)\cap D_+(f_j)),
\end{tikzcd}
\end{center}
其中每一行中的第二、第三个水平箭头分别由
\[
s\mapsto\bigl(s|_{D_+(f_i)}\bigr),\qquad
(s_i)\mapsto\bigl(s_j|_{D_+(f_i)\cap D_+(f_j)}-s_i|_{D_+(f_i)\cap D_+(f_j)}\bigr)
\]
定义。可以验证 $\phi(U)$ 与开覆盖 $\{D_+(f_i)\}_{i\in I}$ 的选取无关。
\end{proof}

\begin{tnote}
$\Gamma_*$ 可看作把 $\Proj S$ 上的层「读回」成分次模的函数：命题 1.4.22 与引理 1.4.23 合起来说明，当 $S$ 由 $S_1$ 作为 $S_0$-代数有限生成时，$M\mapsto M^\sim$ 与 $\cF\mapsto\Gamma_*(\cF)$ 在分次 $S$-模与 $\Proj S$ 的拟凝聚层之间互为（近）逆。注意这里要求 $S$ 由 $S_1$ 有限生成，否则 $\Gamma_*(\cF)^\sim\to\cF$ 未必是同构。[依据：原书 p.86 命题 1.4.22、引理 1.4.23 与 p.82 命题 1.4.19]
\end{tnote}

\begin{lemma}\label{it:1.4.23}
设 $S$ 是环 $A[x_0,\dots,x_m]$，其分次由多项式的次数定义。有 $\Gamma_*(\cO_{\Proj S})\cong S$。
\end{lemma}

\begin{proof}
记 $\cO_{\Proj S}(n)$ 为 $\cO(n)$。因 $D_+(x_i)$（$i=0,\dots,m$）构成 $\Proj S$ 的开覆盖，有正合序列
\[
0\to\cO(n)(\Proj S)\to\prod_{i=0}^m\cO(n)(D_+(x_i))\to\prod_{i,j=0}^m\cO(n)(D_+(x_i)\cap D_+(x_j)),
\]
其中第二个箭头由 $s\mapsto(s|_{D_+(x_i)})$ 定义，第三个箭头由
\[
(s_i)\mapsto\bigl(s_j|_{D_+(x_i)\cap D_+(x_j)}-s_i|_{D_+(x_i)\cap D_+(x_j)}\bigr)
\]
定义。对所有整数 $n$ 取直和，得到正合序列
\[
0\to\Gamma_*(\cO_{\Proj S})\to\prod_{i=0}^mA[x_0,\dots,x_m]_{x_i}\to\prod_{i,j=0}^mA[x_0,\dots,x_m]_{x_ix_j}.
\]
有典范单同态
\[
A[x_0,\dots,x_m]_{x_i}\to A[x_0,\dots,x_m]_{x_ix_j}\to A[x_0,\dots,x_m]_{x_0\cdots x_m}.
\]
经由这些单同态，把 $A[x_0,\dots,x_m]_{x_i}$（$i=0,\dots,m$）与 $A[x_0,\dots,x_m]_{x_0\cdots x_m}$ 的子环等同。上述正合序列说明
\[
\Gamma_*(\cO_{\Proj S})=\bigcap_{i=0}^mA[x_0,\dots,x_m]_{x_i}.
\]
$A[x_0,\dots,x_m]_{x_0\cdots x_m}$ 中的任意元素可唯一写成
\[
x_0^{k_0}\cdots x_m^{k_m}f,
\]
其中 $k_0,\dots,k_m$ 是整数，$f\in A[x_0,\dots,x_m]$ 是不含因子 $x_i$（$i=0,\dots,m$）的多项式。对每个 $i$，$x_0^{k_0}\cdots x_m^{k_m}f\in A[x_0,\dots,x_m]_{x_i}$ 当且仅当对所有 $j\ne i$ 有 $k_j\ge0$。故 $x_0^{k_0}\cdots x_m^{k_m}f\in\bigcap_{i=0}^mA[x_0,\dots,x_m]_{x_i}$ 当且仅当对所有 $j$ 有 $k_j\ge0$。于是 $\Gamma_*(\cO_{\Proj S})=A[x_0,\dots,x_m]$。
\end{proof}

\begin{corollary}\label{it:1.4.24}
\begin{enumerate}[label=(\roman*),leftmargin=2.8em]
  \item 设 $X$ 是 $\bP_A^m=\Proj A[x_0,\dots,x_m]$ 的闭子概形。则存在 $A[x_0,\dots,x_m]$ 的齐次理想 $\fa$，使 $X$ 同构于
  \[
  \Proj\bigl(A[x_0,\dots,x_m]/\fa\bigr).
  \]
  \item $\Spec A$ 上的概形 $X$ 射影，当且仅当存在分次环 $S$，使 $X\cong\Proj S$，其中 $S_0=A$，且 $S$ 由 $S_1$ 作为 $S_0$-代数有限生成。
\end{enumerate}
\end{corollary}

\begin{proof}
(i) 设 $\cI$ 是闭浸入 $X\to\bP_A^m$ 的理想层。由 1.4.22，有 $\cI\cong\Gamma_*(\cI)^\sim$。因
\[
\Gamma_*(\cI)\subset\Gamma_*(\cO_{\bP_A^m})\cong A[x_0,\dots,x_m],
\]
$\Gamma_*(\cI)$ 对应 $A[x_0,\dots,x_m]$ 的齐次理想 $\fa$。则 $X$ 同构于 $\Proj A[x_0,\dots,x_m]/\fa$。

(ii) 留给读者证明「当」的部分。我们证明「仅当」的部分。设 $X$ 在 $\Spec A$ 上射影。则态射 $X\to\Spec A$ 可分解为复合 $X\to\bP_A^m\to\Spec A$，使第一个箭头是闭浸入。由 (i)，$X$ 同构于 $\Proj A[x_0,\dots,x_m]/\fa$，其中 $\fa$ 是某个齐次理想。设 $\fa'=\fa\cap(x_0,\dots,x_m)$。可以证明
\[
\Proj A[x_0,\dots,x_m]/\fa\cong\Proj A[x_0,\dots,x_m]/\fa'.
\]
注意 $S=A[x_0,\dots,x_m]/\fa'$ 具有所需性质。
\end{proof}

对任意概形 $S$，定义 $\bP_S^m=\bP_{\bZ}^m\times S$ 上的\emph{扭层} $\cO_{\bP_S^m}(1)$ 为 $p^*(\cO_{\bP_{\bZ}^m}(1))$，其中 $p\colon\bP_{\bZ}^m\times S\to\bP_{\bZ}^m$ 是投影。当 $S=\Spec A$ 时，它与 $\bP_A^m$ 的扭层一致。一般地，若用仿射开子概形 $U_i=\Spec A_i$（$i\in I$）覆盖 $S$，则 $\cO_{\bP_S^m}(1)$ 由粘合扭层 $\cO_{\bP_{A_i}^m}(1)$（$i\in I$）得到。

设 $X$ 是 $S$-概形。可逆 $\cO_X$-模 $\cL$ 称为在 $S$ 上\emph{极丰沛}\idx{jifengpei}{极丰沛}，若存在 $S$-态射的浸入 $i\colon X\to\bP_S^m$，使 $\cL\cong i^*(\cO_{\bP_S^m}(1))$。

称 $\cO_X$-模 $\cF$ 由整体截面生成，若存在一族截面 $s_i\in\cF(X)$（$i\in I$），使对每个 $P\in X$，$\cF_P$ 作为 $\cO_{X,P}$-模由芽 $s_{iP}$（$i\in I$）生成。这等价于说态射
\[
\bigoplus_{i\in I}\cO_X\to\cF,\quad e_i\mapsto s_i
\]
是满的，其中对每个 $i$，$e_i$ 是 $\bigoplus_{i\in I}\cO_X$ 的第 $i$ 个分量为 $1$、其余分量为 $0$ 的截面。例如，若 $X=\Spec A$，则对每个 $A$-模 $M$，$M$ 中的元素定义 $M^\sim$ 的整体截面，它们生成 $M^\sim$。若 $S$ 是由 $S_1$ 作为 $S_0$-代数生成的分次环，$X=\Proj S$，则 $S_1$ 中的元素定义 $\cO_X(1)$ 的整体截面，它们生成 $\cO_X(1)$。

\begin{proposition}\label{it:1.4.25}[Serre]
设 $X$ 是诺特环 $A$ 上在 $\Spec A$ 上固有的概形，$\cO_X(1)$ 是在 $\Spec A$ 上极丰沛的可逆 $\cO_X$-模，$\cF$ 是凝聚 $\cO_X$-模。则存在整数 $N$，使对每个 $n\ge N$，$\cO_X$-模 $\cF(n)=\cF\otimes_{\cO_X}\cO_X(1)^{\otimes n}$ 由有限多个整体截面生成，即存在满态射
\[
\cO_X^m\to\cF(n)
\]
对某个正整数 $m$。
\end{proposition}

\begin{proof}
设 $i\colon X\to\bP_A^m$ 是使 $i^*(\cO_{\bP_A^m}(1))=\cO_X(1)$ 的浸入。因 $X$ 在 $\Spec A$ 上固有，由 1.3.26 (vi)，$i$ 固有，从而是闭浸入。由下面的 1.4.26，有 $i_*(\cF(n))\cong(i_*\cF)(n)$。此外，对每个 $P\in X$，有 $\cF(n)_P\cong((i_*\cF)(n))_{i(P)}$。为证断言，只需证 $(i_*\cF)(n)$ 对大 $n$ 由有限多个整体截面生成。故可设 $X=\bP_A^m$。对每个 $i$，设 $M_i$ 是有限生成 $A[x_0,\dots,x_m]_{(x_i)}$-模，使 $\cF|_{D_+(x_i)}\cong M_i^\sim$。设 $\{s_{ij}\}$ 是 $M_i$ 的一组有限生成元。则由 1.4.21，存在自然数 $N$，使对任意 $n\ge N$，每个截面 $s_{ij}\otimes x_i^n\in\cF(n)(D_+(x_i))$ 都可延拓为整体截面 $t_{ij}\in\cF(n)(\bP_A^m)$。则 $\{t_{ij}\}$ 生成 $\cF(n)$。
\end{proof}

\begin{tnote}
这是 Serre 在 FAC 中的基本结论（常称「Serre 定理」）：射影概形上的凝聚层扭足够多次后，由有限个整体截面生成。它与「丰沛」概念的对偶表述（$n\ll0$ 时没有非零截面）一起，是射影态射与丰沛层的核心工具。[依据：原书 p.88 命题 1.4.25；J.-P. Serre, Faisceaux algébriques cohérents, Ann. of Math. 61 (1955)]
\end{tnote}

\begin{lemma}\label{it:1.4.26}
设 $f\colon X\to Y$ 是环层空间态射，$\cF$ 是 $\cO_X$-模，$\cG$ 是局部自由 $\cO_Y$-模。则有典范同构
\[
f_*\cF\otimes_{\cO_Y}\cG\cong f_*(\cF\otimes_{\cO_X}f^*\cG).
\]
\end{lemma}

\begin{proof}
设 $\cF'$ 是 $\cO_X$-模。对 $Y$ 的每个开子集 $V$，有典范同态
\[
\cF(f^{-1}(V))\otimes_{\cO_Y(V)}\cF'(f^{-1}(V))\to(\cF\otimes_{\cO_X}\cF')(f^{-1}(V)).
\]
于是得到典范态射
\[
f_*\cF\otimes_{\cO_Y}f_*\cF'\to f_*(\cF\otimes_{\cO_X}\cF').
\]
定义
\[
f_*\cF\otimes_{\cO_Y}\cG\to f_*(\cF\otimes_{\cO_X}f^*\cG)
\]
为复合
\[
f_*\cF\otimes_{\cO_Y}\cG\to f_*\cF\otimes_{\cO_Y}f_*f^*\cG\to f_*(\cF\otimes_{\cO_X}f^*\cG).
\]
为证它是同构，可作局部化并设 $\cG$ 自由。此时断言显然。
\end{proof}

\begin{exercises}
\item \label{ex:1.4.1} 设 $(X,\cO_X)$ 是环层空间，$\cF$ 是 $\cO_X$-模。称 $\cF$ 凝聚，若 $\cF$ 是有限型的，且对每个开子集 $U$ 与每个同态 $u\colon\cO_U^n\to\cF|_U$，$u$ 的核是有限型的。
\begin{enumerate}[label=(\arabic*),leftmargin=2.2em]
  \item 设 $\cO_X$ 凝聚。证明 $\cO_X$-模 $\cF$ 凝聚当且仅当 $\cF$ 有限展示。
  \item 证明对诺特概形，此凝聚定义与本节中的定义一致。
\end{enumerate}

\item \label{ex:1.4.2} 设 $X$ 与 $Y$ 是 $S$-概形，$f,g\colon X\to Y$ 是 $S$-态射，$U$ 是 $X$ 的稠开子集，使 $f|_U=g|_U$ 作为概形态射。设 $Y$ 在 $S$ 上分离。则 $f=g$ 作为拓扑空间之间的映射。再设 $X$ 既约。则 $f=g$ 作为概形态射。

\item \label{ex:1.4.3} 设 $S$ 是由 $S_1$ 作为 $S_0$-代数生成的分次环。
\begin{enumerate}[label=(\arabic*),leftmargin=2.2em]
  \item 证明对每个 $f\in S_1$ 与每个分次 $S$-模 $M$，有 $S_{(f)}$-模同构
  \[
  \bigoplus_{n=-\infty}^\infty M(n)_{(f)}\xrightarrow{\sim}M_f.
  \]
  特别地，取 $M=S$，得 $S_{(f)}$-模同构
  \[
  \bigoplus_{n=-\infty}^\infty S(n)_{(f)}\xrightarrow{\sim}S_f.
  \]
  证明有 $S_f$-模同构
  \[
  M_{(f)}\otimes_{S_{(f)}}S_f\xrightarrow{\sim}M_f.
  \]
  \item 设 $M$ 与 $N$ 是分次 $S$-模。构造同构
  \[
  M_{(f)}\otimes_{S_{(f)}}N_{(f)}\xrightarrow{\sim}(M\otimes_SN)_{(f)},
  \]
  使在取张量积 $\otimes_{S_{(f)}}S_f$ 并与 (1) 中的同构复合后，得到典范同构
  \[
  M_f\otimes_{S_f}N_f\cong(M\otimes_SN)_f.
  \]
  \item 设 $\fp\in\Proj S$，$T$ 是 $S-\fp$ 中齐次元素的集合。证明有 $S_{(\fp)}$-模同构
  \[
  \bigoplus_{n=-\infty}^\infty M(n)_{(\fp)}\xrightarrow{\sim}T^{-1}M.
  \]
  特别地，取 $M=S$，得 $S_{(\fp)}$-模同构
  \[
  \bigoplus_{n=-\infty}^\infty S(n)_{(\fp)}\xrightarrow{\sim}T^{-1}S.
  \]
  证明有 $T^{-1}S$-模同构
  \[
  M_{(\fp)}\otimes_{S_{(\fp)}}T^{-1}S\xrightarrow{\sim}T^{-1}M.
  \]
  \item 设 $M$ 与 $N$ 是分次 $S$-模。构造同构
  \[
  M_{(\fp)}\otimes_{S_{(\fp)}}N_{(\fp)}\xrightarrow{\sim}(M\otimes_SN)_{(\fp)},
  \]
  使在取张量积 $\otimes_{S_{(\fp)}}T^{-1}S$ 并与 (3) 中的同构复合后，得到典范同构
  \[
  T^{-1}M\otimes_{T^{-1}S}T^{-1}N\xrightarrow{\sim}T^{-1}(M\otimes_SN).
  \]
\end{enumerate}

\item \label{ex:1.4.4} 设 $X$ 是概形，$Y$ 是 $X$ 的闭子概形，$\cI$ 是 $Y$ 的理想层。定义 $X$ 沿 $Y$ 的\emph{爆破}\idx{baopo}{爆破} $\tilde X$ 为 $X$ 上的概形 $\Proj\bigl(\bigoplus_{n=0}^\infty\cI^n\bigr)$，其中 $\Proj\bigl(\bigoplus_{n=0}^\infty\cI^n\bigr)$ 由粘合 $X$ 的仿射开子概形 $U$ 上的 $\Proj\bigl(\bigoplus_{n=0}^\infty\cI^n|_U\bigr)$ 得到。
\begin{enumerate}[label=(\arabic*),leftmargin=2.2em]
  \item 设 $\pi\colon\tilde X\to X$ 是典范态射。证明 $\pi$ 诱导同构 $\pi^{-1}(X-Y)\xrightarrow{\sim}X-Y$。
  \item 设 $f\colon X'\to X$ 是闭浸入，$Y'=X'\times_XY$ 是由基变换 $Y\to X$ 得到的 $X'$ 的闭子概形。证明 $X'$ 沿 $Y'$ 的爆破 $\tilde X'$ 同构于复合
  \[
  X'-Y'\xrightarrow{f}X-Y\xrightarrow{\pi^{-1}}\pi^{-1}(X-Y)\to\tilde X
  \]
  的概形论像。
\end{enumerate}

\item \label{ex:1.4.5} 设 $S$ 是概形，$\cF$ 是有限秩局部自由 $\cO_S$-模，$\cF^\vee=\mathcal{H}om_{\cO_S}(\cF,\cO_S)$ 是 $\cF$ 的对偶，$\operatorname{Sym}^n(\cF^\vee)$ 是 $\cF^\vee$ 的 $n$ 次对称积，且
\[
\cB=\bigoplus_{n=0}^\infty\operatorname{Sym}^n(\cF^\vee).
\]
证明 $\cB$ 是拟凝聚 $\cO_S$-代数。设 $f\colon\Spec\cB\to S$ 是使 $f_*\cO_{\Spec\cB}\cong\cB$ 的仿射态射。考虑层 $\cF'$，使对 $S$ 的每个开子集 $U$，$\cF'(U)$ 是 $S$-态射 $U\to\Spec\cB$ 的集合。证明 $\cF\cong\cF'$。$S$-概形 $\Spec\cB$ 称为与 $\cF$ 相伴的\emph{向量丛}。

\item \label{ex:1.4.6} 设 $X$ 是整概形，$K$ 是 $X$ 的函数域，$L$ 是含 $K$ 的域。对 $X$ 的每个仿射开子集 $U=\Spec A$，可把 $K$ 与 $A$ 的分式域等同。设 $\tilde A$ 是 $A$ 在 $L$ 中的整闭包，即 $L$ 中满足形如
\[
x^n+a_1x^{n-1}+\cdots+a_n=0
\]
的代数方程（系数 $a_i$ 在 $A$ 中）的所有元素 $x$ 组成的环。对任意仿射开子集 $V=\Spec B$，$V\subset U$，有典范同态 $\tilde A\to\tilde B$。证明存在拟凝聚 $\cO_X$-代数 $\tilde{\cO}$，使对每个仿射开子集 $U=\Spec A$ 有 $\tilde{\cO}(U)\cong\tilde A$，且对每对仿射开子集 $V\subset U$，限制 $\tilde{\cO}(U)\to\tilde{\cO}(V)$ 可等同于典范同态 $\tilde A\to\tilde B$。$X$-概形 $\tilde X=\Spec\tilde{\cO}$ 称为 $X$ 在 $L$ 中的\emph{整闭包}。$X$ 在函数域 $K$ 中的整闭包称为 $X$ 的\emph{正规化}。若 $X$ 同构于其正规化，则称 $X$ 是\emph{正规}的。证明 $X$ 正规当且仅当对每个 $P\in X$，$\cO_{X,P}$ 整闭。
\end{exercises}
